% ══════════════════════════════════════════════════════════════════════════════
%  INTRODUCTION GÉNÉRALE
% ══════════════════════════════════════════════════════════════════════════════

Depuis quelques années, le numérique a considérablement transformé les méthodes d'enseignement et d'apprentissage. Au Maroc, cette évolution est portée par la Vision Stratégique 2015-2030 et le programme Maroc Digital 2030, qui accordent une place importante à la modernisation du système éducatif. Malgré ces avancées, les étudiants de l'enseignement supérieur continuent de recevoir l'essentiel de leurs cours sous format PDF, un format qui reste difficile à exploiter sans outils interactifs adaptés.

C'est dans ce contexte qu'a été conçu \textbf{EduAI}, une plateforme SaaS (\textit{Software as a Service}) qui s'appuie sur l'Intelligence Artificielle, et plus précisément sur le Traitement Automatique du Langage Naturel (NLP) et les grands modèles de langage (LLM), pour permettre aux étudiants de transformer leurs cours PDF en supports d'apprentissage interactifs.

Ce rapport présente le travail réalisé dans le cadre du Projet de Fin d'Études de la filière Bachelor en Ingénierie Informatique en Intelligence Artificielle et Sciences des Données, à l'École Supérieure de Technologie d'Essaouira (Université Cadi Ayyad). L'idée du projet est partie d'un constat simple : beaucoup d'étudiants peinent à réviser efficacement à partir de cours PDF souvent longs et denses.

\textbf{EduAI} est un projet personnel que j'ai réalisé seul, sous l'encadrement de \textbf{Pr. El Mahdi ERRAJI}. En parallèle de mes études, je travaille comme Développeur Full-stack et Responsable des Opérations chez \textbf{2Pi eLearning SARL}, qui a bien voulu mettre à disposition un espace sur son infrastructure Portainer pour héberger la plateforme.

\vspace{0.5cm}
Le présent rapport est organisé en cinq chapitres :

\begin{enumerate}[label=\textbf{Chapitre~\arabic* :}]
  \item \textbf{Contexte général et problématique} — Présentation de l'organisme d'accueil, analyse du contexte éducatif marocain, formulation de la problématique et étude critique de l'existant.
  \item \textbf{Analyse et spécification des besoins} — Identification des acteurs, spécification des besoins fonctionnels et non-fonctionnels, et modélisation UML des cas d'utilisation.
  \item \textbf{Conception} — Architecture globale du système, conception des pipelines NLP/IA, modélisation de la base de données et conception de l'API REST.
  \item \textbf{Réalisation} — Environnement technique, choix technologiques, présentation des modules implémentés et validation par les tests.
  \item \textbf{Déploiement, modèle économique et perspectives} — Stratégie de déploiement conteneurisé, modèle SaaS freemium, projections financières et axes d'évolution.
\end{enumerate}

Les diagrammes UML et les captures d'écran des interfaces sont regroupés en annexes.
